% eksperymentalne tłumaczenie na włoski

%

\chapter{Geometrie nieeuklidesowe} % lang-pl
\chapter{Geometrie non euclidee} % lang-it


Tekst sekcji geometrie nieeuklidesowe. % lang-pl
Testo della sezione sulle geometrie non euclidee. % lang-it


\section{Geometria Łobaczewskiego} % lang-pl
\section{Geometria di Lobachevskij} % lang-it


Eves: sekcja 7.2, 7.3. % lang-pl
Eves: sezioni 7.2, 7.3. % lang-it

Limit triangles. % lang-pl
Triangoli limite. % lang-it

Istnieje naturalna miara długości odcinków. % lang-pl
Esiste una misura naturale della lunghezza dei segmenti. % lang-it

Czworokąt Saccheriego. % lang-pl
Quadrilatero di Saccheri. % lang-it

Suma kątów trójkąta jest $< \pi$. % lang-pl
La somma degli angoli di un triangolo è $< \pi$. % lang-it

Pole trójkąta, defekty. % lang-pl
Area del triangolo, difetti angolari. % lang-it

Punkty idealne, ultraidealne, linie hiperrównoległe. % lang-pl
Punti ideali, ultraideali e rette iperparallele. % lang-it


\section{Dobre, stare} % lang-pl
\section{I classici} % lang-it

\begin{proposition}[aksjomat Arystotelesa] % lang-pl
	Aksjomat Arystotelesa. % lang-pl
\begin{proposition}[assioma di Aristotele] % lang-it
	
	Assioma di Aristotele. % lang-it
\index{aksjomat!Arystotelesa}%
\end{proposition}

\begin{proposition}[lemat Proklusa] % lang-pl
	Prosta nie może przecinać tylko jednej z dwóch prostych równoległych. % lang-pl
\begin{proposition}[lemma di Proclo] % lang-it
	
	Una retta non può intersecare soltanto una di due rette parallele. % lang-it
\index{lemat!Proklusa}%
\end{proposition}

\begin{proposition}[aksjomat Claviusa] % lang-pl
	Aksjomat Claviusa. % lang-pl
\begin{proposition}[assioma di Clavio] % lang-it
	
	Assioma di Clavio. % lang-it
\index{aksjomat!Claviusa}%
\end{proposition}

\begin{proposition}[aksjomat Clairauta] % lang-pl
	Aksjomat Clairauta. % lang-pl
\begin{proposition}[assioma di Clairaut] % lang-it
	
	Assioma di Clairaut. % lang-it
\index{aksjomat!Clairauta}%
\end{proposition}

\begin{proposition}[aksjomat Simsona] % lang-pl
	Aksjomat Simsona. % lang-pl
\begin{proposition}[assioma di Simson] % lang-it
	
	Assioma di Simson. % lang-it
\index{aksjomat!Simsona}%
\end{proposition}

\begin{proposition}[aksjomat Playfaira] % lang-pl
	Aksjomat Playfaira. % lang-pl
\begin{proposition}[assioma di Playfair] % lang-it
	
	Assioma di Playfair. % lang-it
\index{aksjomat!Playfaira}%
\end{proposition}

Treść aksjomatu Playfaira poda szkocki matematyk John Playfair w podręczniku \emph{Elements of Geometry} z 1795 roku. % lang-pl

L'enunciato dell'assioma di Playfair fu formulato dal matematico scozzese John Playfair nel manuale \emph{Elements of Geometry} del 1795. % lang-it
% % https://en.wikipedia.org/wiki/Playfair%27s_axiom
\index[persons]{Playfair, John}%

\begin{proposition}[aksjomat Wallisa] % lang-pl
	Aksjomat Wallisa. % lang-pl
\begin{proposition}[assioma di Wallis] % lang-it
	
	Assioma di Wallis. % lang-it
\index{aksjomat!Wallisa}%
\end{proposition}

\begin{proposition}[aksjomat Bolyi] % lang-pl
	Aksjomat Bolyi. % lang-pl
\begin{proposition}[assioma di Bolyai] % lang-it
	
	Assioma di Bolyai. % lang-it
\index{aksjomat!Bolyi}%
\end{proposition}

\begin{definition}[czworokąt Saccheriego] % lang-pl
	Czworokąt Saccheriego. % lang-pl
\begin{definition}[quadrilatero di Saccheri] % lang-it
	
	Quadrilatero di Saccheri. % lang-it
% TODO: https://en.wikipedia.org/wiki/Lambert_quadrilateral
\index{czworokąt!Saccheriego}
\end{definition}

% TODO: https://www.deltami.edu.pl/2016/09/inne-swiaty-inne-geometrie/

Giovanni Girolamo Saccheri będzie uczniem Tommasa Cevy, brata Giovanniego, autora twierdzenia Cevy (w tej książce z numerem \ref{twierdzenie_cevy}). % lang-pl

Giovanni Girolamo Saccheri sarà allievo di Tommaso Ceva, fratello di Giovanni, autore del teorema di Ceva (in questo libro numerato come \ref{twierdzenie_cevy}). % lang-it

\begin{proposition}[aksjomat Legendre'a] % lang-pl
	Aksjomat Legendre'a. % lang-pl
\begin{proposition}[assioma di Legendre] % lang-it
	
	Assioma di Legendre. % lang-it
\index{aksjomat!Legendre'a}%
\end{proposition}

\begin{definition}[model Poincarego] % lang-pl
	Model Poincarego. % lang-pl
\begin{definition}[modello di Poincaré] % lang-it
	
	Modello di Poincaré. % lang-it
\index{model Poincarego}%
\end{definition}

\begin{definition}[geometria hiperboliczna] % lang-pl
	Geometria hiperboliczna. % lang-pl
\begin{definition}[geometria iperbolica] % lang-it
	
	Geometria iperbolica. % lang-it
\index{geometria!hiperboliczna}
\end{definition}

O geometrii hiperbolicznej wspomina się w $\Delta_{24}^7$. % lang-pl
La geometria iperbolica è menzionata in $\Delta_{24}^7$. % lang-it
\begin{proposition} % lang-it
	Se la geometria euclidea del piano è coerente, allora anche la geometria di Lobachevskij del piano è coerente. % lang-it
\end{proposition} % lang-it


\begin{proposition} % lang-pl
	Jeśli geometria euklidesowa (na płaszczyźnie) jest niesprzeczna, to geometria Łobaczewskiego (na płaszczyźnie) też jest niesprzeczna. % lang-pl
\end{proposition} % lang-pl
Eves \cite[s. 347-353]{eves1_1972}. % lang-pl
Eves \cite[p. 347-353]{eves1_1972}. % lang-it


W 1733 roku Giovanni Girolamo Saccheri spróbuje zbadać, co by było gdyby piąty postulat Euklidesa okazał się fałszywy. % lang-pl
Nel 1733 Giovanni Girolamo Saccheri cercherà di capire che cosa accadrebbe se il quinto postulato di Euclide fosse falso. % lang-it

Jego prace zapoczątkują systematyczne rozważania nad alternatywami dla geometrii euklidesowej. % lang-pl
I suoi lavori daranno inizio allo studio sistematico delle alternative alla geometria euclidea. % lang-it

W 1829 roku János Bolyai, Carl Friedrich Gauss oraz (niezależnie od nich) Nikołaj Łobaczewski stworzą geometrię hiperboliczną, czyli pierwszą w pełni rozwiniętą geometrię nieeuklidesową, w której piąty postulat Euklidesa nie obowiązuje. % lang-pl
Nel 1829 János Bolyai, Carl Friedrich Gauss e, indipendentemente da loro, Nikolaj Lobachevskij costruiranno la geometria iperbolica, la prima geometria non euclidea pienamente sviluppata, nella quale il quinto postulato di Euclide non vale. % lang-it


Kolejnym osiągnięciem będzie model Beltramiego-Kleina dla geometrii hiperbolicznej, w którym pracujemy z punktami wewnątrz otwartego koła. % lang-pl
Un ulteriore progresso sarà il modello di Beltrami-Klein della geometria iperbolica, nel quale si lavora con i punti all'interno di un disco aperto. % lang-it

Proste w tym modelu to cięciwy łączące punkty z brzegu koła. % lang-pl
Le rette di questo modello sono le corde che uniscono punti del bordo del disco. % lang-it

Wspomni o nim Eugenio Beltrami w 1868 roku, najpierw dla $n = 2$, a potem dla dowolnej liczby wymiarów. % lang-pl
Eugenio Beltrami lo presenterà nel 1868, dapprima per $n=2$ e successivamente per un numero arbitrario di dimensioni. % lang-it
% memoirs % lang-it

% memoirs % lang-pl
Wcześniej, bo w 1859 roku, Arthur Cayley użyje dwustosunkowej definicji kąta (od Laguerre'a) by pokazać, że geometria euklidesowa może być zdefiniowana przy pomocy geometrii rzutowej. % lang-pl
In precedenza, nel 1859, Arthur Cayley utilizzerà la definizione di angolo tramite rapporto anarmonico (dovuta a Laguerre) per mostrare che la geometria euclidea può essere definita mediante la geometria proiettiva. % lang-it

Cayley zapozna się z pracami Beltramiego i Cayleya; w 1870 roku da wykład na seminarium u Weierstrassa, który zakończy pytaniem: czy między pomysłami Cayleya oraz Łobaczewskiegi istnieje powiązanie? % lang-pl
Felix Klein studierà i lavori di Beltrami e Cayley; nel 1870 terrà un seminario presso Weierstrass, concludendolo con la domanda: esiste una relazione tra le idee di Cayley e quelle di Lobachevskij? % lang-it

% na podstawie https://en.wikipedia.org/wiki/Beltrami–Klein_model
% gdzie pada
% and these essays proved the equiconsistency of hyperbolic geometry with ordinary Euclidean geometry

% model Kleina - Delta 2012-06

% Sferyczna osobno?
% https://en.wikipedia.org/wiki/Lexell's_theorem
\begin{proposition}[twierdzenie Lexella] % lang-pl
	Wierzchołki trójkątów sferycznych o ustalonej podstawie $AB$ i polu powierzchni leżą na okręgu małym (powstałym przez przecięcie sfery płaszczyzną) przechodzącym przez punkty antypodalne do wierzchołków $A$, $B$ z podstawy. % lang-pl
\begin{proposition}[teorema di Lexell] % lang-it
	
	I vertici dei triangoli sferici con base fissata $AB$ e area fissata giacciono su un cerchio minore (ottenuto intersecando la sfera con un piano) passante per i punti antipodali dei vertici $A$ e $B$ della base. % lang-it
\end{proposition}
% https://en.wikipedia.org/wiki/Spherical_trigonometry => Girard

Anders Johan Lexell napisze około 1777 roku pracę na temat tego twierdzenia, razem z dwoma dowodami: trygonometrycznym i geometrycznym. % lang-pl

Inni podadzą inne dowody: % lang-pl
Anders Johan Lexell scriverà intorno al 1777 un lavoro dedicato a questo teorema, contenente due dimostrazioni: una trigonometrica e una geometrica. % lang-it

Leonhard Euler w~1778, % lang-pl
Altri ne forniranno ulteriori dimostrazioni: % lang-it

Adrien-Marie Legendre w~1800, % lang-pl
Leonhard Euler nel 1778, % lang-it

Jakob Steiner w~1827, % lang-pl
Adrien-Marie Legendre nel 1800, % lang-it

Carl Friedrich Gauss w~1841, % lang-pl
Jakob Steiner nel 1827, % lang-it

Paul Serret w~1855, % lang-pl
Carl Friedrich Gauss nel 1841, % lang-it

Joseph-Émile Barbier w~1864. % lang-pl
Paul Serret nel 1855, % lang-it

Joseph-Émile Barbier nel 1864. % lang-it

% https://www.deltami.edu.pl/2012/06/w-rozumowaniach-byl-blad/
% https://www.deltami.edu.pl/2018/08/geometria-bolyaia-lobaczewskiego-co-to-jest-i-jak-ja-poznawac/

%